\section{C. The normalized periods, quantitative bounds, and the boundary}

\subsection{C.1. A theta expression for the missing period}
We use the odd Jacobi function
\[
 \vartheta_1(v,\tau)=2\sum_{n\geq0}(-1)^n
 e^{\pi i\tau(n+1/2)^2}\sin((2n+1)\pi v)
\]
and the Dedekind function
\(\eta_D(\tau)=e^{\pi i\tau/12}\prod_{n\geq1}(1-e^{2\pi in\tau})\).
The subscript distinguishes this function from the invariant integral class.
We choose the cusp lift reached from \(\tau=iT\), \(T>1\), on which
\(\mu=1/2-iy(T)\), \(0<y(T)<T/2\). Put
\(a=5+2\sqrt6\) and \(\rho=e^{\pi i/3}\).

\begin{proposition}[C.1.1; Proved]\label{st:C.1.1}
For the period section identified in Statement B.1.3, the particular solution
with \((\beta^{\rm part}+\tau)(p_0)=0\) is
\begin{equation}\label{eq:cd-beta}
 \beta^{\rm part}=\frac{6i}{\pi}\operatorname{Log}
 \frac{\vartheta_1(\mu,\tau)}{2\sqrt3\,\eta_D(\tau)}.
\end{equation}
Here the logarithm is continued on the simply connected uniformizing
half-plane from its cusp expansion. Its special values are
\[
 \mu_0=\frac12-\frac{i\log a}{2\pi},\qquad
 \mu(z_1)=\frac{2-\rho}{3},\qquad \mu(z_2)=\frac{1-i}{2},
\]
\[
 \beta^{\rm part}(z_2)=i\left(\frac32-\frac3\pi\log6\right),
 \qquad \operatorname{Re}\beta^{\rm part}(z_1)=-\frac12.
\]
In particular \(e^{2\pi i\mu_0}=-a\).
\end{proposition}
\begin{proof}
Statement B.1.3 identifies the elliptic point represented by \(\mu\) with
one of the sections \((0,\pm\sqrt2t^2(t-1))\). The uniformization map
\((v,\tau)\mapsto(f^2\wp_\tau(v),f^3\wp'_\tau(v)/2)\) therefore gives
\(\wp_\tau(\mu)=0\). No zero of \(\wp\) is a lattice point. Thus the numerator
of \eqref{eq:cd-beta} is nonzero; the denominator is nonzero as well.
A holomorphic logarithm exists on the simply connected cover.

For completeness the cusp value follows directly from the normally
convergent Fourier expansion, rather than from a quoted Eichler--Zagier
constant:
\begin{equation}\label{eq:cd-wp}
 \frac{\wp_\tau(v)}{\pi^2}
 =\csc^2(\pi v)-\frac13+
 8\sum_{n\geq1}\frac{nq^n}{1-q^n}(1-\cos(2\pi nv)),
 \qquad q=e^{2\pi i\tau}.
\end{equation}
This identity follows by differentiating the logarithm of the Jacobi
product twice and using
\(E_2=1-24\sum nq^n/(1-q^n)\). Boundedness of the chosen elliptic logarithm
at the cusp permits passage to the limit. We obtain
\(\sin^2(\pi\mu_0)=3\), equivalently
\(u^2+10u+1=0\), \(u=e^{2\pi i\mu_0}\).
On the chosen imaginary-axis branch \(\operatorname{Im}\mu<0\), so \(u=-a\).
One can see the branch without a numerical convention: on the rectangular
elliptic curve with \(\tau=iT\), \(T>1\), \(\wp\) is real on the segment
\(1/2-iy\), \(0\leq y\leq T/2\). Its endpoint values are the largest and
middle real roots of its cubic, respectively positive and negative.
Its derivative has no zero in the open segment, since the zeros of
\(\wp'\) are the three nonzero half-periods. There is exactly one zero
there; this gives the indicated continuation. Consequently
\(2\sin(\pi\mu_0)=2\sqrt3\).

Put \(R=\vartheta_1/\eta_D\). Poisson summation for the defining Gaussian
series and its index shift give
\[
 \begin{aligned}
 R(v,\tau+1)&=e^{\pi i/6}R(v,\tau),\\
 R(v/\tau,-1/\tau)&=-i e^{\pi iv^2/\tau}R(v,\tau),\\
 R(1-v,\tau)&=R(v,\tau),\qquad R(v+1,\tau)=-R(v,\tau).
 \end{aligned}
\]
Inserting the actual transformed arguments gives
\[
 \begin{aligned}
 R(\mu\circ g_1,\tau\circ g_1)
 &=e^{-\pi i/3}e^{\pi i(1-\mu)^2/\tau}R(\mu,\tau),\\
 R(\mu\circ g_2,\tau\circ g_2)
 &=i e^{\pi i\mu^2/\tau}R(\mu,\tau).
 \end{aligned}
\]
Multiplication of the logarithms by \(6i/\pi\) gives exactly
\(2-6(1-\mu)^2/\tau\) and \(-3-6\mu^2/\tau\), modulo constants in
\(12\mathbb Z\). These constants vanish: each defect is constant on the
connected cover, and evaluation at the corresponding fixed point makes
both its two sides zero. The fixed-point equations themselves give
\(\mu(z_1)=(2-\rho)/3\) and \(\mu(z_2)=(1-i)/2\).
The Jacobi product has cusp expansion
\(R(\mu,\tau)=2q^{1/12}\sin(\pi\mu)(1+O(q))\).
Equation \eqref{eq:cd-beta} consequently gives
\(\beta^{\rm part}+\tau\to0\). The difference from the specified particular
solution is an invariant holomorphic function extending across all three
points of \(B\), hence a constant, and the cusp normalization makes it zero.

At \(\tau=i\), the half-period translation formula gives
\(\vartheta_1((1-i)/2,i)=e^{\pi/4}\vartheta_3(0,i)\).
The identities \(2\eta_D^3=\vartheta_2\vartheta_3\vartheta_4\),
\(\vartheta_2(0,i)=\vartheta_4(0,i)=2^{-1/4}\vartheta_3(0,i)\)
give \(R=e^{\pi/4}\sqrt2\), proving the displayed value of \(\beta\).
At \(\rho\), \(q\) is negative real, \(\mu=1/2-i/(2\sqrt3)\), and the
product for \(R\) has phase \(\pi/12\), with all its remaining factors
positive real. This gives \(\operatorname{Re}\beta=-1/2\).
\end{proof}

\begin{proposition}[C.1.2; Proved]\label{st:C.1.2}
The following decimal evaluations of the exact expressions in C.1.1 agree
in two independent evaluations through the displayed twenty decimal places:
\[
\begin{array}{c|l}
\mu_0&0.5-0.36485183191106784792\,i\\
\mu(z_1)&0.5-0.28867513459481288225\,i\\
\mu(z_2)&0.5-0.5\,i\\
\beta^{\rm part}(z_1)&-0.5-1.27101206057537346287\,i\\
\beta^{\rm part}(z_2)&-0.21100425815613411418\,i
\end{array}
\]
These are evaluations of explicit convergent expressions, with working
precision ninety decimal digits. The formula \eqref{eq:cd-beta}, together
with \(j(\tau)=1728t\) and \(\wp_\tau(\mu)=0\), determines the periods
throughout a fundamental domain.
\end{proposition}
\begin{proof}
The executable \texttt{cd\_numerics.py} evaluates \(\mu_0\) first by its
algebraic logarithm and then by solving
\(\pi^2\csc^2(\pi v)-\pi^2/3=0\). For \(\beta\), its first calculation uses
the Jacobi theta series and the eta product; its second uses
\[
 R(v,\tau)=2e^{\pi i\tau/6}\sin(\pi v)
 \prod_{n\geq1}(1-2q^n\cos(2\pi v)+q^{2n}).
\]
The product is truncated after 120 factors. At the two specified points
\(|q|\leq e^{-\pi\sqrt3}\) and \(|\operatorname{Im}v|\leq1/2\), so the
omitted logarithm is bounded by
\(\sum_{n>120}(2|q|^n\cosh\pi+|q|^{2n})/(1-2|q|^n\cosh\pi-|q|^{2n})\),
which is less than \(10^{-280}\). The measured discrepancies of the
ninety-digit calculations are below \(10^{-88}\). We do not promote floating-point rounding to an interval certificate.
A further sixty-digit interval computation, \texttt{cd\_special\_intervals.py},
evaluates the real positive product at $\rho$, including its forty-factor
tail, and the exact logarithmic expression at $i$. It encloses both
imaginary parts within intervals narrower than $3\cdot10^{-59}$ and
therefore certifies the displayed twenty decimal places. It also proves
$-0.2<\operatorname{Im}(\kappa_1^{\rm part}-\kappa_2^{\rm part})<-0.1$.
The exact formulas, product-tail bounds, independent algorithms, and
full outputs accompany the decimals. The special values of \(\mu\) are also checked directly in
the fixed-point equations, an exact algebraic calculation.
\end{proof}

\subsection{C.2. Bounds for the sharp parameter disc}
\begin{proposition}[C.2.1; Proved]\label{st:C.2.1}
For the normalization of C.1.1, the exact admissibility maximum is finite
and satisfies
\[
 -\frac3\pi\log6\ \leq M<2.
\]
More precisely the proof below gives \(M<1.635\). Consequently the sharp
radius \(R_w=e^{-\pi M}\) satisfies
\(e^{-2\pi}<R_w\leq216\), and for the reference linearization with
\(u=w^2\), \(e^{-4\pi}<R_u\leq46656\).
\end{proposition}
\begin{proof}
Equation \eqref{eq:cd-beta} gives
\begin{equation}\label{eq:cd-norm}
 D_{\rm part}=\frac6\pi\log\frac{\mathcal N(\mu,\tau)}{2\sqrt3},
 \qquad
 \mathcal N(v,\tau)=e^{-\pi(\operatorname{Im}v)^2/\operatorname{Im}\tau}
 |R(v,\tau)|.
\end{equation}
The theta transformation formulas prove that \(\mathcal N\) is unchanged
by integral period translations of \(v\), sign change of \(v\), and the
simultaneous modular change of \(v,\tau\). We may therefore use the
comparison isomorphism \(\mathbb C/(\mathbb Z+\tau\mathbb Z)\to
\mathbb C/(\mathbb Z+\tau'\mathbb Z)\), \([v]\mapsto[v/(c\tau+d)]\), to
place \(\tau\) in the standard modular domain and then choose
\(|\operatorname{Im}v|\leq T/2\), \(T=\operatorname{Im}\tau\geq\sqrt3/2\).
Both operations preserve the equation \(\wp(v)=0\) and \eqref{eq:cd-norm}.

For \(\sqrt3/2\leq T\leq2\), the theta series and
\(\prod(1-r^n)\geq\exp(-r/(1-r)^2)\) imply
\[
 \mathcal N\leq
 2e^{\pi/6}\left(1+\frac{e^{-\pi\sqrt3/2}}
 {1-e^{-3\pi\sqrt3/2}}\right)
 \exp\frac{e^{-\pi\sqrt3}}{(1-e^{-\pi\sqrt3})^2}
 <3.615.
\]
Indeed the absolute value of the \(n\)-th sine term, including the
Gaussian factor, is at most
\(2\exp(-\pi T(n+1/2-|\operatorname{Im}v|/T)^2)\), which is at most
\(2e^{-\pi Tn^2}\); and \(n^2\geq1+3(n-1)\) for \(n\geq1\).

For \(T\geq2\), a zero of \(\wp\) in this period strip has
\(|\operatorname{Im}v|<1\). To prove this, assume its absolute value is at
least one and apply \eqref{eq:cd-wp}. With \(r=e^{-4\pi}\), the absolute
value of the terms other than \(-1/3\) is at most
\[
 \frac1{\sinh^2\pi}+8\left\{
 \frac r{(1-r)^3}+\frac1{2(1-r)}
 \left(\frac{e^{-2\pi}}{(1-e^{-2\pi})^2}
 +\frac{e^{-6\pi}}{(1-e^{-6\pi})^2}\right)\right\}<0.016<\frac13.
\]
Here \(|\sin\pi v|\geq\sinh\pi\), and
\(|q|^n\cosh(2\pi n\operatorname{Im}v)\leq
(e^{-\pi Tn}+e^{-3\pi Tn})/2\). This contradicts \(\wp(v)=0\).
The product expression for \(R\) now gives the uniform bound
\[
 \mathcal N\leq2\cosh\pi\,e^{-\pi/3}
 \exp\left(\frac{2\cosh(2\pi)r}{1-r}+\frac{r^2}{1-r^2}\right)
 <8.151.
\]
Insertion in \eqref{eq:cd-norm} proves the stated upper bound. The script
\texttt{cd\_bounds.py} independently evaluates these monotone bounds by
fifty-digit interval arithmetic: the larger resulting bound for \(D\) is
contained, after outward rounding for display, in
\[
 [1.63421874605891012363,\;1.63421874605891012365].
\]
The output file retains the full fifty-digit endpoints.
This encloses the coarse upper bound; the exact maximum is proved in C.2.3.

At the cusp \(D_{\rm part}=-T+o(1)\), so it tends to minus infinity;
continuity on the rest of the compact orbifold base gives attainment of
its maximum. At \(z_2\), C.1.1 gives
\(D_{\rm part}=-(3/\pi)\log6\), proving the lower bound. Addition of \(c\)
adds \(\operatorname{Im}c\) to \(D\); a zero is forced by continuity
between a maximum and the cusp whenever \(\operatorname{Im}c\geq-M\).
For \(\operatorname{Im}c<-M\), the real determinant \(TD\) is everywhere
nonzero, proving the exact asserted form of the admissible disc.
\end{proof}

\begin{proposition}[C.2.2; Proved]\label{st:C.2.2}
The point \(t=1\) is a strict local maximum of \(D_{\rm part}\).
A double-precision calculation on 5151 points of the half modular domain
\(0\leq\operatorname{Re}\tau\leq1/2\), \(|\tau|\geq1\),
\(\operatorname{Im}\tau\leq3\), has largest sampled value at \(\tau=i\).
This finite computation does not establish the global maximum.
\end{proposition}
\begin{proof}
We give a quantitative neighbourhood, which will also specify the local
region needed for a global interval cover. Put
\[
 h=v-(1-\tau)/2,\quad x=\tau-i,\quad
 G(h,\tau)=\wp_\tau((1-\tau)/2+h),\quad
 H(h,\tau)=\log\frac{\vartheta_3(h,\tau)}{\eta_D(\tau)}.
\]
The half-period theta identity gives exactly
$\log\mathcal N=\operatorname{Re}H-\pi(\operatorname{Im}h)^2/T$.
On the complex polydisc $|h|\leq1/5$, $|x|\leq1/10$, the theta series
has $|\vartheta_3-1|<1$, so the logarithm is holomorphic. There are no
poles of $G$ there: for every integer lattice point its imaginary
difference from $(1-\tau)/2+h$ has absolute value at least
$0.9/2-0.2>0$. Both functions are even in $h$.

The following uniform bounds use only positive Gaussian sums:
\[
 |G|<20,\qquad |H|<0.6.
\]
For explicit verification put
$A_k=\sum_{n\geq1}n^k e^{-0.9\pi n^2}\cosh(0.4\pi n)$,
$a_0=2A_0$, $a_1=4\pi A_1$, $a_2=8\pi^2A_2$, and $r=e^{-1.8\pi}$.
The bounds used are
\[
 \begin{aligned}
 |G|&\leq\frac{a_2}{1-a_0}+\frac{a_1^2}{(1-a_0)^2}
 +\frac{\pi^2}3\left(1+\frac{24r}{(1-r)^3}\right)<18.373,\\
 |H|&\leq-\log(1-a_0)+\frac{1.1\pi}{12}+
 \frac r{(1-r)^2}<0.547.
 \end{aligned}
\]
The derivative identity behind the first bound is
$G=-\partial_h^2\log\vartheta_3-\pi^2E_2/3$.
The script \texttt{cd\_local\_bound.py} evaluates these bounds by sixty-digit
interval arithmetic. It truncates the positive sums at $n=10$ and bounds
the remaining terms by geometric sums using
$n^2\geq121+23(n-11)$; the polynomial factors $n^k$ are summed explicitly.

At $(0,i)$ the theta heat equation gives
\[
 G=0,\quad G_{h^2}=-g_2(i)/4,\quad G_x=id,
 \qquad H=\log\sqrt2,\quad H_x=0,\quad H_{h^2}=-\pi/2.
\]
Subscripts $h^2$ denote power-series coefficients. The first two follow
also from $\wp''=6\wp^2-g_2/2$. Here $d$ is the real positive derivative coefficient of
$e_2(\tau)=(\pi^2/3)(\vartheta_2(0,\tau)^4-
\vartheta_4(0,\tau)^4)$ at $i$. Differentiating the convergent theta
series, with explicit Gaussian tails, gives $10.03<d<10.04$.
The value $H_x=0$ follows by differentiating the modular $S$-invariance
of $\vartheta_3(0,\tau)/\eta_D(\tau)$ at $i$. Differentiating
$\eta_D(-1/\tau)=\sqrt{-i\tau}\eta_D(\tau)$ gives $E_2(i)=3/\pi$,
which with $\wp((1-i)/2,i)=0$ gives $H_{h^2}=-\pi/2$.
The positive Eisenstein series for $g_2$, with an explicit tail in the
same script, gives $47.26<g_2(i)/4<47.30$.

Set $\epsilon=3/20000$, $r_h=3/100$, $s=\epsilon/(1/10)=0.0015$,
and $r_2=(r_h/(1/5))^2=0.0225$.
Cauchy's estimates for the even power series of $G$ give, on
$|h|=r_h$, $x=0$,
\[
 |G(h,i)|\geq
 \left(47.26-\frac{500r_2}{1-r_2}\right)r_h^2>0.03217.
\]
For $|x|\leq\epsilon$, the change of $G$ on this circle is at most
$20s/((1-s)(1-r_2))<0.03074$. Rouche's theorem therefore places both
zeros of $G(\cdot,\tau)$ in $|h|<r_h$. They represent the two zeros of
$\wp$ modulo its lattice, since an elliptic function with one double
pole has two zeros with multiplicity.

Write $G=e_2(\tau)+h^2Q(h,\tau)$. The same Cauchy bounds give
\[
 |Q+g_2(i)/4|\leq
 \frac{500(r_2+s/(1-s))}{1-r_2},
\quad
 |e_2(\tau)-idx|\leq
 \frac{2000|x|^2}{1-s}.
\]
At any zero they imply $|x|\leq7|h|^2$ and $|h|^2<1/20000$;
the interval output records bounds $6.124$ and $0.00004434$,
respectively. For $H$, its stated coefficients and Cauchy's estimates
now bound the remainder, together with the change of $1/T$ in the
Gaussian, by $5000|h|^4$. For example a direct upper bound for its
coefficient is
\[
 \frac{0.6(70^2+25\cdot70+25^2)}
 {(1-s)(1-25/20000)}+\frac{7\pi}{1-\epsilon}<5000.
\]
Consequently, for $0<|x|\leq3/20000$,
\[
 \log\mathcal N-\log\sqrt2
 \leq-\frac\pi2|h|^2+5000|h|^4<0.
\]
At $h=0$ the equation $G=0$ forces $x=0$ in this disc by the lower
bound for $e_2$. This proves strict local maximality, with an explicit
radius, and covers the square
$|\operatorname{Re}\tau|\leq10^{-4}$,
$|\operatorname{Im}\tau-1|\leq10^{-4}$.

The grid calculation uses thirty Fourier terms for $\wp$, complex
Newton iteration, and the theta product. Five sample values were
independently recomputed with eighty Fourier terms and the theta series
at ninety-digit precision; discrepancies are below $7\cdot10^{-16}$.
The source, precise grid rule, sampled maximum, and comparisons are in
\texttt{cd\_numerics.py} and its output.
\end{proof}

\begin{theorem}[C.2.3; Proved, with interval certificate]\label{st:C.2.3}
For the cusp normalization $\beta^{\rm part}+\tau\to0$,
\[
 M=-\frac3\pi\log6=-1.71100425815613411418\ldots,
 \qquad \operatorname*{arg\,max}_{B\setminus\{p_0\}}D_{\rm part}=\{1\}.
\]
The exact parameter discs are
\[
 0<|w|<216,
 \qquad 0<|u|<46656
\]
for the reference linearization $u=w^2$; an arbitrary reference constant
$\lambda_0$ multiplies the latter radius by $|\lambda_0|$.
The proof combines an analytic cusp bound, the quantitative local proof
C.2.2, and a completed outward-rounded interval cover of the remaining
compact domain.
\end{theorem}
\begin{proof}
By \eqref{eq:cd-norm} it is enough to prove
$\mathcal N(v,\tau)\leq\sqrt2$ at every zero of $\wp_\tau$, with equality
only in the modular orbit of $i$. Modular transformations, period
translations and sign change preserve this norm. Complex conjugation
also preserves it, through
$(v,\tau)\mapsto(\bar v,-\bar\tau)$, as follows by conjugating the
normally convergent theta and eta series. It therefore suffices to work
on $0\leq\operatorname{Re}\tau\leq1/2$, $|\tau|\geq1$.

For $T\geq2$, the zero-strip argument in C.2.1 improves to
$|\operatorname{Im}v|<b=43/100$. Replace $1/\sinh^2\pi$ in its displayed
error bound by $1/\sinh^2(\pi b)$; the resulting upper bound is less
than $0.316<1/3$, as certified independently in
\texttt{cd\_special\_intervals.py}. Retain the Gaussian in the product
estimate. The function
$\cosh(\pi y)\exp(-\pi y^2/T)$ increases on $0\leq y\leq b$, $T\geq2$:
its logarithmic derivative is $\pi\tanh(\pi y)-2\pi y/T\geq0$, since
$\tanh(\pi y)\geq y$ there. The latter follows from concavity and its
endpoint values. With $r=e^{-4\pi}$ we get
\[
 \mathcal N\leq
 2\cosh(\pi b)e^{-\pi/3-\pi b^2/2}
 \exp\left(\frac{2\cosh(2\pi b)r}{1-r}+
                  \frac{r^2}{1-r^2}\right)
 <1.082<\sqrt2.
\]
The bound decreases for $T\geq2$: the logarithmic derivative of the
exponential $e^{-\pi T/6-\pi b^2/T}$ is negative and the product
majorant decreases with $|q|$. This excludes the whole cusp region,
without a numerical truncation of the base.

We now specify the compact computation as a proof by enclosure.
Its rational starting rectangle is
$\mathcal R=[0,1/2]+i[17/20,2]$, which contains the compact half modular
domain. For every rational rectangle $Q$ retained in this domain, the
program selects a decimal complex centre $v_0$, a decimal nonzero
preconditioner $A$, and a closed complex rectangle $Z$ around $v_0$.
Floating-point Newton iteration only selects these numbers. The proof
uses the subsequent interval tests
\[
 v_0-A\wp(v_0,Q)+(1-A\partial_v\wp(Z,Q))(Z-v_0)
       \subset\operatorname{int}Z,
 \qquad
 \|1-A\partial_v\wp(Z,Q)\|_\infty<1.
\]
Here the norm of multiplication by a complex number is bounded by the
sum of the absolute real and imaginary parts. All displayed sets are
outward-rounded intervals. The mean-value integral along a segment in
$Z$ proves that, for every $\tau\in Q$, the map
$v\mapsto v-A\wp_\tau(v)$ maps $Z$ into itself and is a contraction.
Banach's theorem supplies a genuine zero in $Z$. Both zeros of $\wp$
are related by sign and lattice translations, so their norms agree;
this one enclosure therefore controls the actual period section.

For reproducibility the Fourier remainders are explicit. With
$q=e^{2\pi i\tau}$, $p=q e^{2\pi iv}$, $m=q e^{-2\pi iv}$, let
$r\geq|q|$ and $s\geq\max(|p|,|m|)$ be the interval upper endpoints.
The code requires $s<1$. For $N=24$ put
\[
 S_1(r)=r^{N+1}\left(\frac{N+1}{1-r}+\frac r{(1-r)^2}\right),
\]
\[
 S_2(r)=r^{N+1}\left(\frac{(N+1)^2}{1-r}+
 \frac{2(N+1)r}{(1-r)^2}+\frac{r(1+r)}{(1-r)^3}\right).
\]
The errors in $\wp$ and $\partial_v\wp$ are bounded respectively by
$8\pi^2(S_1(r)+S_1(s))/(1-r)$ and $16\pi^3S_2(s)/(1-r)$.
These follow by absolute summation of \eqref{eq:cd-wp} and its derivative.
They are added as rectangles enclosing the corresponding complex discs.
For the norm, use $h=v-(1-\tau)/2$ and the even theta series
$\vartheta_3=1+2\sum_{n\geq1}e^{\pi i n^2\tau}\cos(2\pi nh)$.
With $T_*\leq\operatorname{Im}\tau$ and $H_*\geq|\operatorname{Im}h|$,
its omitted terms after $n=10$ have absolute value at most
\[
 \frac{2e^{-121\pi T_*+22\pi H_*}}
 {1-e^{-23\pi T_*+2\pi H_*}},
\]
and the denominator is tested positive. The logarithm of the eta
product tail after $n=17$ is bounded by
$r^{18}/((1-r)(1-r^{18}))$. The resulting enclosure of
$e^{-2\pi(\operatorname{Im}h)^2/T}|\vartheta_3/\eta_D|^2$
is required to have upper endpoint strictly below two.

The completed computation \texttt{f\_interval\_parallel.py}, using
\texttt{f\_interval\_cert.py}, runs these tests at thirty decimal digits
and records every rational box. It accepted 19701 norm boxes. A further
36 boxes lie inside the local square of C.2.2 and are covered by that
analytic proof. Another 174 boxes satisfy
$x_{\rm upper}^2+y_{\rm upper}^2<1$ exactly and lie outside the modular
domain. There are zero pending or undecided boxes. The independent
program \texttt{f\_interval\_verify.py} reconstructs the entire binary
partition from its leaves, checks all reported strict interval bounds,
and verifies the exclusions by rational arithmetic. It consumes all
19911 leaves, with exact total area $23/40$ and maximum depth 38.
The largest reported squared-norm upper bound is the exact decimal
rational
\[
 \frac{9999999982782673130994674616307871}
 {5000000000000000000000000000000000}<2.
\]
The JSON certificate, complete outputs, and independently runnable
source programs are included. Thus the compact domain is fully covered;
the exploratory grid is not used as a substitute for any enclosure.
The local proof gives equality only at $\tau=i$, and the cusp and other
compact boxes give strict inequality. This proves the global bound and
its equality set.

Finally C.1.1 evaluates $D$ at $i$ exactly. Since $j(\tau)=1728t$, its
modular orbit maps precisely to $t=1$. Exponentiation gives
$e^{-\pi M}=e^{3\log6}=216$ and $e^{-2\pi M}=6^6=46656$.
The two independent theta calculations and the sixty-digit logarithmic
interval computation in C.1.2 verify the numerical value of $M$ as well
as the exact evaluation; the integer radii are checked against interval
exponentiation in \texttt{cd\_special\_intervals.py}.
\end{proof}

\subsection{C.3. The rank-three period quotient}
\begin{proposition}[C.3.1; Proved]\label{st:C.3.1}
At every point with \(D=0\), \(T>0\), the real-linear map
\(\Pi:\mathbb R^4\to\mathbb C^2\) has rank three. Write
\(\mu=a_0\tau+b_0\), \(a_0,b_0\in\mathbb R\), and put
\[
 k=(1,-6a_0,-6b_0,6a_0\operatorname{Re}\mu-\operatorname{Re}\beta).
\]
Its real kernel is \(\mathbb R k\). The subgroup \(\Pi\mathbb Z^4\) is
discrete exactly when \(k\in\mathbb Q^4\). In that case it is a
rank-three lattice in a real hyperplane and its quotient in \(\mathbb C^2\)
is a noncompact complex Lie group. In every other case the topological
quotient \(\mathbb C^2/\Pi\mathbb Z^4\) is non-Hausdorff.
\end{proposition}
\begin{proof}
The real matrix, with output order
\((\operatorname{Re}\zeta_1,\operatorname{Im}\zeta_1,
\operatorname{Re}\zeta_2,\operatorname{Im}\zeta_2)\), is
\[
 \begin{pmatrix}6\operatorname{Re}\mu&\operatorname{Re}\tau&1&0\\
 6\operatorname{Im}\mu&T&0&0\\
 \operatorname{Re}\beta&\operatorname{Re}\mu&0&1\\
 \operatorname{Im}\beta&\operatorname{Im}\mu&0&0\end{pmatrix}.
\]
Its last three columns are real-linearly independent because \(T>0\).
Its determinant is \(TD\), so at zero determinant the rank is exactly
three. Substitution gives the displayed kernel, including its signs.
If the image subgroup is discrete, it is a lattice of rank at most three
in its real span. The homomorphism \(\mathbb Z^4\to\Pi\mathbb Z^4\)
therefore has a nonzero kernel, forcing the displayed kernel vector with
first coordinate one to be rational. Conversely a rational line has a
primitive integral generator extendable to a basis of \(\mathbb Z^4\).
The remaining three basis images are a real basis of the image hyperplane
and hence generate a discrete lattice. Its quotient leaves one unbounded
real direction. Finally a non-discrete finitely generated subgroup of
this real hyperplane cannot be closed: a closed subgroup is a product of
a real vector group and a discrete group, and a nonzero real vector group
is uncountable. A quotient topological group is Hausdorff exactly when its
identity subgroup is closed. This proves the last assertion.
\end{proof}

\begin{proposition}[C.3.2; Proved]\label{st:C.3.2}
The unique degenerate base point at the actual admissibility boundary is
\(t=1\). The boundary constants
\(c=x+i(3/\pi)\log6\) give
\[
 \ker_{\mathbb R}\Pi=\mathbb R(1,3,-3,-x-3/2).
\]
For rational $x+3/2=m/n$ in lowest terms, the quotient is a
finite holomorphic quotient of $E_i\times\mathbb C^*$, with a covering
of degree four. For irrational \(x+3/2\), the closure of the period subgroup is generated
by the real line \(\mathbb R(0,1)\) and the two periods
\((i,(1-i)/2)\), \((1,0)\). Its Hausdorff quotient is a real manifold
homeomorphic to \(\mathbb R\times(S^1)^2\).
\end{proposition}
\begin{proof}
C.1.1 gives \(\operatorname{Im}\beta=3/2\) at this value of \(c\).
The vector \(\widehat\gamma+3\widehat u-3\widehat w\) maps to
\((0,x+3/2)\), while \(\widehat\delta\) maps to \((0,1)\).
For $x+3/2=m/n$ in lowest terms, change coordinates by
$(\zeta_1,\zeta_2)\mapsto(\zeta_1,\zeta_2+\zeta_1/2)$.
The vectors $2\widehat u-\widehat\delta$ and
$2\widehat w-\widehat\delta$ then have periods $(2i,0)$ and $(2,0)$.
Bezout's identity applied to $\widehat\gamma+3\widehat u-3\widehat w$
and $\widehat\delta$ gives the period $(0,1/n)$.
These three periods generate a sublattice of index four in the rank-three
image: its two remaining generators are represented by
$(i,1/2)$ and $(1,1/2)$, each of order two modulo this sublattice.
The corresponding quotient is $E_i\times(\mathbb C/(n^{-1}\mathbb Z))$,
and $[z]\mapsto e^{2\pi inz}$ identifies its second factor with
$\mathbb C^*$. Lattice inclusion gives the claimed degree-four map.
For irrational $x+3/2$, its additive subgroup with $1$ is dense in the
indicated real line. The other two displayed periods are independent
modulo that line. This proves the closure and the quotient description.
By C.2.3 this is the only degenerate torus cover at the admissibility
boundary; every other finite-base period matrix retains real rank four.
We have described the torus cover of the order-four filling, including
the actual quotient and its topology.
\end{proof}

\begin{remark}[C.3.3; Open]\label{st:C.3.3}
The remaining compactification problem at the actual boundary is now
localized at the order-four fibre over $t=1$. Its smallest undecided
instance is a compact complex-space replacement of the rank-three
quotient in C.3.2, over a neighbourhood of one boundary parameter,
with a proof that it extends the original cyclic filling and its overlap
with the ordinary torus family. Both rational and irrational values of
$x+3/2$ are included in the exact quotient description. The naive
quotients cannot be central fibres of a proper complex-space family:
for rational $x+3/2$ they are noncompact, and for irrational values
they are non-Hausdorff. C.3.4 supplies a further obstruction to preserving
a holomorphic Weierstrass Albanese model in this parameter.
These prove failure of the stated specific replacements. A compact
modification that changes that model has not been constructed or ruled
out. The requested general complex-space boundary filling remains
unfinished, with its unique base-fibre location and exact local period
obstruction determined.
\end{remark}

\begin{proposition}[C.3.4; Proved]\label{st:C.3.4}
On its half-plane of definition, the $j$-invariant of
$\operatorname{Alb}(S_2)$ has no meromorphic continuation across any
point of the line
$\operatorname{Im}c=(3/\pi)\log6$. In particular the rank-three quotient
at $t=1$ cannot be replaced by a Weierstrass family of elliptic Albanese
curves whose coefficients are holomorphic in the same parameter across
that line.
\end{proposition}
\begin{proof}
The explicit Albanese lattice in A.4 is
$\mathbb Z+(\kappa_2/2)\mathbb Z$, with
$\kappa_2=c+3/2-i(3/\pi)\log6$ by C.1.1. A positive-imaginary modulus is
$z=-\kappa_2/2$, so the invariant is $j(z)$. At a reduced rational
$r=a/b$ on the real boundary, choose an integral determinant-one modular
matrix with bottom row $(b,-a)$. Its value at $r+iy$ has imaginary part
$1/(b^2y)$. The convergent Fourier expansion $j=q^{-1}+744+O(q)$ therefore
gives growth of order $\exp(2\pi/(b^2y))$ as $y\downarrow0$.
A meromorphic function at $r$ is bounded by a fixed power of $1/y$,
so continuation there is impossible. A meromorphic continuation near an
irrational real point would also continue the same function near some
rational real point in that neighbourhood, which has just been excluded.
The affine change from $c$ to $z$ transfers this obstruction to the
stated line. A Weierstrass family with holomorphic coefficients has a
meromorphic $j$-invariant, expressed as the quotient of its cubic
invariants, so it cannot supply the indicated continuation. This
obstruction concerns the stated Albanese construction; it does not
silently assume that every possible compact threefold filling preserves
an Albanese model of its special surface.
\end{proof}

\subsection{C.4. Classification of the non-normal cusp surface}
\begin{theorem}[C.4.1; Proved]\label{st:C.4.1}
For the normalization of C.1.1, the three opposite-side identification
constants of \(W_c\), in the characters specified below, are
\[
 (g_1,g_2,g_3)=(-a^{-1},a^6,-a^{-7}e^{-2\pi ic}).
\]
Two such compact complex surfaces are biholomorphic exactly when
\(c'-c\in\mathbb Z\). The vertical \(\mathbb C^*\)-action inherited from
translation in the \(\widehat\delta\)-direction fixes exactly the double
curve represented by the pair of edges of direction \(e_2\).
\end{theorem}
\begin{proof}
Keep the lattice of toric cocharacters \(\mathbb Z^2\), with
\(e_1=\widehat w\), \(e_2=\widehat\delta\), and cusp logarithm
\(s=(2\pi i)^{-1}\log t_c\). The cusp period matrix is
\([sB_0+C(t_c)\mid I]\), where
\[
 B_0=\begin{pmatrix}0&1\\-1&0\end{pmatrix},\qquad
 C(0)=\begin{pmatrix}6\mu_0&h_0\\b_0-h_0&\mu_0\end{pmatrix},
 \quad b_0=c.
\]
For \(v=e_1,e_2,e_1-e_2\), the gluing map from the side of direction
\(v\) to its opposite is the restriction of the toric action with lattice
label \(B_0v\). A character row \(r\) with \(rv=0\) pulls back by this
map to itself multiplied by
\(\exp(2\pi i\,rC(0)B_0v)\): the remaining power of \(t_c\) is
\(rB_0^2v=-rv=0\). Taking respectively
\(r=(0,1),(1,0),(1,1)\) gives exponents
\(-\mu_0,6\mu_0,-7\mu_0-c\). This proves the three constants, including
the cancellation of \(h_0\). These are morphisms of the whole projective
boundary lines, since multiplication by a nonzero constant extends over
both endpoints. The restrictions agree at both triple points, as the
ambient toric action supplies the same vertex identifications.

A biholomorphism of the two surfaces lifts uniquely to a biholomorphism
of their normalizations: integral closure of the reduced local rings is
functorial under an isomorphism. It preserves the inverse image of the
singular locus, the hexagon on \(\mathrm{dP}_6\). Its permutation of the
six components preserves their cyclic incidence, hence lies in the
order-twelve dihedral group. After composition with the corresponding
toric lattice automorphism, it fixes each boundary component. Each torus
character has its divisor supported on those components, so its pullback
divided by itself has divisor zero on the compact normalization and is
constant. Such an automorphism is therefore multiplication by an element
of \((\mathbb C^*)^2\). These torus multipliers cancel in the opposite-side
comparison because the same character is used on both sides.

The remaining six effective actions on the triple, with the central
lattice automorphism \(-I\) acting trivially, are described by the equations
\(g'_{k_j}{}^{\epsilon_j}=g_j\), where
\[
 ((k_1,\epsilon_1),(k_2,\epsilon_2),(k_3,\epsilon_3))
 \in\left\{
\begin{array}{lll}
((1,+),(2,+),(3,+)),&((1,-),(3,+),(2,+)),\\
((2,-),(3,-),(1,+)),&((2,+),(1,+),(3,-)),\\
((3,-),(2,-),(1,-)),&((3,+),(1,-),(2,-)).
\end{array}\right.
\]
This finite calculation is reproduced by \texttt{cd\_exact.py}, enumerating
all integral automorphisms of the six rays and transporting the characters.
Every nonidentity action requires one of
\(g_1^{-1}=g_1\), \(g_2^{-1}=g_2\), or equality of a positive real number
\(g_2^{\pm1}\) with a negative real number \(g_1^{\pm1}\).
All fail for \(a>1\). Therefore \(g'_3=g_3\), precisely
\(e^{-2\pi ic'}=e^{-2\pi ic}\). Conversely equality of the three constants
gives the identity gluing morphism on the normalization and hence a
biholomorphism of the quotients. This proves the unmarked assertion.

For the action of the subtorus with cocharacter \(e_2\), a one-dimensional
toric orbit on the normalization is pointwise fixed exactly when its
stabilizer line contains \(e_2\). Among the six rays this occurs exactly
for \(\pm e_2\). Their images form one double curve; all zero-dimensional
orbits map to its two endpoints. No dense two-dimensional orbit is fixed.
The quotient action has the same fixed set: a continuous orbit cannot
jump between the finitely many branches of a normalization fibre.
\end{proof}

\section{D. Exceptional fibres, Neron--Severi classes, and compact subspaces}
\subsection{D.1. Elliptic quotients and the torsion locus}
\begin{proposition}[D.1.1; Proved]\label{st:D.1.1}
The specified linear projection \(\operatorname{pr}_1:\mathbb C^2\to
\mathbb C\) induces a holomorphic group homomorphism
\(F_t\to E_t=\mathbb C/(\mathbb Z\tau+\mathbb Z)\) exactly on
\(6\mu=A\tau+B\), \(A,B\in\mathbb Z\). It is then surjective with
connected kernel
\(\mathbb C/(\mathbb Z+(A\mu-\beta)\mathbb Z)\).
On the smooth base the two values of \(t\) are precisely
\[
 t_\pm=\frac{64(47\pm27\sqrt3)}{11}.
\]
\end{proposition}
\begin{proof}
The images of the four lattice generators under the linear projection
are \(6\mu,\tau,1,0\). This proves descent in both directions.
The integral basis change replacing \(\widehat\gamma\) with
\(\widehat\gamma-A\widehat u-B\widehat w\) changes its period to
\((0,\beta-A\mu)\); the remaining vertical generator is \((0,1)\).
The other two periods map onto the entire target lattice, so the kernel
is connected. Its second modulus has imaginary part
\(A\operatorname{Im}\mu-\operatorname{Im}\beta=-D>0\).

The elliptic model may be scaled by the explicit isomorphism
\((X,Y)\mapsto(2X,2\sqrt2Y)\) to
\[
 Y^2=X^3-\frac34t^3(t-1)X+\frac14t^4(t-1)^2,
 \qquad P=(0,t^2(t-1)/2).
\]
Its sixth division polynomial at \(P\) is
\[
 \psi_6(P)=\frac9{65536}t^{24}(t-1)^9
 (11t^2-6016t+8192).
\]
For a smooth elliptic fibre, \(\psi_n(P)=0\) is exactly the condition
\(nP=O\), since the multiplication rational map has these and only these
poles in its affine coordinate. The two roots displayed above are simple
and avoid \(0,1\). We computed the same equation independently by the
chord-and-tangent addition algorithm, obtaining
\[
 \begin{aligned}
 x(6P)&=\frac{(89t^2+4736t-4096)Q(t)}
 {1296(11t^2-6016t+8192)^2},\\
 Q(t)&=881t^4+4745344t^3-12079104t^2\\
 &\hspace{2em}+24641536t-16777216.
 \end{aligned}
\]
The division-polynomial recurrence and the addition algorithm are in
separate supplied programs; all polynomial operations use exact rational
arithmetic. Thus the denominator has no unverified numerical root or
cancellation. The decimal values \(1.36510950191\) and
\(545.543981407\) are only convenient approximations to the exact roots.
\end{proof}

\begin{proposition}[D.1.2; Known, with derivation]\label{st:D.1.2}
The locus on \(B^\circ\) where \(P(t)\) is torsion is countable and dense.
Every fibre over this locus is an abelian surface. The divisor on the
smooth base cut out by \(6kP=O\) has degree \(3k^2-1\), including local
intersection multiplicity, for every positive integer \(k\).
\end{proposition}
\begin{proof}
On a simply connected open subset write uniquely
\(\mu=a_0\tau+b_0\), \(a_0,b_0\) real. The section is torsion precisely
when \((a_0,b_0)\in\mathbb Q^2\). In a local complex parameter \(z\), direct
differentiation gives
\[
 \det\frac{\partial(a_0,b_0)}{\partial(\operatorname{Re}z,
 \operatorname{Im}z)}
 =-\frac{|\mu'(z)-a_0\tau'(z)|^2}{\operatorname{Im}\tau}.
\]
This determinant vanishes on no open set. On an open set of vanishing,
the two real differential equations imply \(da_0=db_0=0\), hence
\(\mu=a_0\tau+b_0\) with constant real coefficients. Analytic continuation
would give this identity on the cover. The cusp law
\((\tau,\mu)\circ g_0=(\tau-1,\mu)\) then forces \(a_0=0\), and the
\(g_2\)-law makes a constant \(\mu\) impossible. Thus every open subset
contains a point where this real map is a local diffeomorphism, and its
image contains rational pairs. This proves density. For each integer
triple \((n,A,B)\), \(n\ne0\), the function \(n\mu-A\tau-B\) is
holomorphic and not identically zero by the same argument. Its zeros are
discrete; a countable union proves countability.

For \(N\mu=A\tau+B\), \(A,B,N\in\mathbb Z\), \(N>0\), use the linear
coordinate change \(\zeta'_2=\zeta_2-(A/N)\zeta_1\).
The four integral vectors
\[
 N\widehat u-B\widehat\delta,\quad
 N\widehat w+A\widehat\delta,\quad
 N\widehat\gamma-6A\widehat u-6B\widehat w,\quad
 \widehat\delta
\]
span a sublattice of index \(N^3\). Its first two periods lie on
\(\zeta'_2=0\), and its last two lie on \(\zeta_1=0\).
Both intersections are genuine rank-two lattices because \(D<0\).
Their product quotient is therefore a product of elliptic curves, and
inclusion of this sublattice induces a finite surjective holomorphic
homomorphism onto \(F_t\). A positive integral form on the product can
be summed over its finite translation kernel and multiplied by the
index to descend an integral positive form to \(F_t\). The Riemann
polarization criterion proves projectivity.

For the degree statement we use the height computation actually read in
\cite[Lemma 2.1 and proof]{E26}: \(\langle P,P\rangle=1/6\),
\(\chi(\mathcal O_S)=1\), the component groups at the two reducible fibres
have orders three and two, and the local corrections for \(P\) are
\(4/3\) and \(1/2\). The application to \(6kP\) is as follows.
The corrected divisor class of a section is its class minus \(O\) and
a fibre multiple, plus the unique rational combination of nonidentity
fibre components that makes it perpendicular to those components.
On restriction to the generic fibre this assignment is additive because
the divisor class \([Q]-[O]\) is the group law in \(\operatorname{Pic}^0\).
The perpendicular correction makes its self-intersection the negative of
the height, so the height is quadratic. The element \(6kP\) is in the
identity component at both reducible fibres; its correction terms vanish.
Hence
\[
 6k^2=(6k)^2\langle P,P\rangle
 =2+2((6kP)\mathbin{\cdot}O),
\]
which gives the formula. There are no missing intersection points at the
three special fibres. The computed rational functions for \(6P\) show
that it has finite affine coordinates at \(0,1\), so it is a nonzero point
of each additive identity component; multiplication by \(k\) remains
nonzero in \(\mathbb C\). At infinity its multiplicative coordinate is
\((-a)^{6k}=a^{6k}\ne1\). Thus all intersections counted by the height
occur on \(B^\circ\). This derives the cited height input in precisely
the form needed here; the first five degrees and the correction matrices
are checked exactly in \texttt{cd\_torsion\_independent.py}.
\end{proof}

\begin{theorem}[D.1.3; Refuted]\label{st:D.1.3}
The assertion in Lead 6 that \emph{existence of a surjection onto the
specified \(E_t\) with elliptic kernel} is equivalent to \(6P(t)=O\)
is false. It remains true for the specified projection in D.1.1.
\end{theorem}
\begin{proof}
Set \(t=32/5\). The exact fourth division polynomial is
\[
 \psi_4(P)=-\frac1{32}t^{10}(t-1)^4(5t-32),
\]
whereas \(\psi_2(P)=t^2(t-1)\ne0\). Thus \(P\) has exact order four.
On the rectangular branch write \(\tau=iT\), \(T>1\), with
\(j(iT)=1728\cdot32/5\). The chosen zero has real part \(1/2\), negative
imaginary part of absolute value less than \(T/2\), and order four;
therefore \(\mu=(2-\tau)/4\). Choose the concrete parameter
\[
 c=\left(\frac98-\frac{101}{4}\right)\tau-\frac94
       -\beta^{\rm part}(t).
\]
It is globally admissible. Indeed the positive-real Jacobi product at this
\((\mu,\tau)\) gives
\(\operatorname{Im}\beta^{\rm part}\geq-T-(3/\pi)\log3\).
Consequently
\(\operatorname{Im}c\leq-185T/8+(3/\pi)\log3<-22<-M\), by C.2.1.
This inequality proves admissibility without using the numerical estimate
\(c=-2.25-34.02109724218209472308\,i\).

The holomorphic linear map \(L:\mathbb C^2\to\mathbb C\),
\(L(\zeta)=3\zeta_1+4\zeta_2\), takes the four periods to
\[
 -101\tau,\qquad2\tau+2,\qquad3,\qquad4.
\]
These generate exactly \(\mathbb Z\tau+\mathbb Z\): the last two generate
\(1\), and the first two then generate \(\tau\), since
\(\gcd(101,2)=1\). Thus \(L\) descends to a surjective holomorphic group
homomorphism \(F_{32/5}\to E_{32/5}\) with connected one-dimensional
kernel. The kernel is compact as a closed subset of a compact torus,
hence an elliptic curve. Nevertheless \(6P=2P\ne O\).
Every equality, the division-polynomial noncancellation, and the four
lattice images are checked by exact arithmetic in the supplied code.
\end{proof}

\begin{proposition}[D.1.4; Proved]\label{st:D.1.4}
For every smooth fibre and every admissible \(c\), its Neron--Severi lattice
is the following explicit subgroup of \(\mathbb Z^6\):
\begin{equation}\label{eq:cd-NS}
 \left\{E:\ E_{12}-E_{13}\tau+(6E_{23}-E_{14})\mu+
 E_{24}\beta+E_{34}(6\mu^2-\tau\beta)=0\right\}.
\end{equation}
The complex orientation has intersection form
\(E^2=-2\operatorname{Pf}(E)\), where
\(\operatorname{Pf}(E)=E_{12}E_{34}-E_{13}E_{24}+E_{14}E_{23}\).
In particular the invariant class \(\eta=6\gamma\wedge\delta+u\wedge w\)
has \(\int_{F_t}\eta^2=-12\). Its Hermitian form is indefinite.
The fibres with lattice exactly \(\mathbb Z\eta\) have no compact curves
and algebraic dimension zero.
\end{proposition}
\begin{proof}
On the universal cover the holomorphic two-form is the wedge of the two
rows of \(\Pi\). Its six coefficients in the integral coframe are
\(6\mu^2-\tau\beta,-\beta,6\mu,-\mu,\tau,1\).
An integral real two-form is of type \((1,1)\) exactly when its wedge
with this form vanishes: its \((0,2)\) part is then zero, and reality
also annihilates its \((2,0)\) part. Expanding the wedge gives
\eqref{eq:cd-NS}. The exponential sequence identifies these integral
\((1,1)\) classes with Chern classes of holomorphic line bundles, proving
the lattice assertion.

The real period matrix in C.3.1 has determinant \(TD<0\). Thus the
integral orientation \((\widehat\gamma,\widehat u,\widehat w,
\widehat\delta)\) is opposite to the complex orientation. Algebraically
\(E\wedge E=2\operatorname{Pf}(E)\gamma\wedge u\wedge w\wedge\delta\),
which proves the sign in the displayed intersection formula. In
particular the expression ``\(\eta^2=12\)'' in AC26, Lemma 9.4 is correct
as its displayed coefficient of \(\gamma\wedge u\wedge w\wedge\delta\),
and has the opposite sign as a complex-oriented numerical
self-intersection. We retain the coefficient identity and record the
orientation comparison explicitly.

For the precise Hermitian convention $H_E$ defined in D.1.6, direct
substitution of the four periods gives $H_\eta=-2G^{-1}$, where
\[
 G=\begin{pmatrix}2T&2U\\2U&\operatorname{Im}\beta/3\end{pmatrix},
 \qquad \det G=2TD/3<0.
\]
Thus $H_\eta$ has one positive and one negative eigenvalue, and every
nonzero multiple remains indefinite. The integration
current of an effective divisor, averaged under translations of the
compact torus, is a nonzero semipositive constant \((1,1)\)-form in its
Chern class. It cannot be an indefinite multiple of \(\eta\).
Thus no curve exists when \(\operatorname{NS}=\mathbb Z\eta\).
A nonconstant meromorphic function would have a nonempty polar divisor
(after subtracting a constant if necessary), so the algebraic dimension
is zero.
\end{proof}

\begin{proposition}[D.1.5; Proved]\label{st:D.1.5}
For a fixed \(c\), all but countably many smooth fibres have
\(\operatorname{NS}=\mathbb Z\eta\). Conversely, for every fixed smooth
\(t\), there are admissible parameters tending to the cusp of the
parameter disc for which \(F_t\) is an abelian surface, irrespective of
the order of \(P(t)\).
\end{proposition}
\begin{proof}
We verify the required independence of five holomorphic functions.
Suppose \(A+B\tau+C\mu+D\beta+E(6\mu^2-\tau\beta)=0\).
Apply \(g_0\), using \((\tau,\mu,\beta)\circ g_0=(\tau-1,\mu,\beta+1)\),
and subtract. The result is \(-B+D+E(\beta-\tau+1)=0\).
A second cusp difference gives \(2E=0\), hence \(E=0,D=B\).
Applying \(g_2\) to the remaining relation and multiplying the difference
by \(-\tau\) gives
\[
 6B\mu^2+C(\tau-1)\mu+B(\tau^2+3\tau+1)-C\tau=0.
\]
Its cusp difference gives \(C\mu+2B(\tau+1)-C=0\), whose next cusp
difference forces \(B=0\). Then \(C\mu=C\); the \(g_2\)-law forces
\(C=0\), and finally \(A=D=0\). Thus an integral class outside
\(\mathbb Z\eta\) has a nonidentically-zero holomorphic left side in
\eqref{eq:cd-NS}; its zeros are discrete. Countability of \(\mathbb Z^6\)
proves the first assertion, with explicit verification of its reduction.

For the second assertion put
\(\beta_n=-3n^2\tau+12n\mu\),
\(c_n=\beta_n-\beta^{\rm part}(t)\). Its imaginary part tends to
minus infinity, so the already proved finite bound on \(M\) makes every
sufficiently large \(n\) admissible. The integral alternating matrix
with nonzero upper entries
\[
 E_{13}=-3n^2,\quad E_{14}=6n,\quad E_{23}=-n,\quad E_{24}=1
\]
satisfies \eqref{eq:cd-NS} exactly at \(\beta_n\), and
\(\operatorname{Pf}(E)=-3n^2<0\). Its complex-oriented square is positive.
A nondegenerate Hermitian form in complex dimension two with positive
determinant is definite. Choosing its sign gives a positive integral
Riemann form, so \(F_t\) is projective. This supplies an explicit family
of projective fibres even at any chosen non-torsion value of the section.
\end{proof}

\begin{theorem}[D.1.6; Known, with derivation]\label{st:D.1.6}
Fix any admissible $c$. For each integral alternating matrix $E$, let
$Z_E(c)$ be the zero set on the marking cover of the explicit holomorphic
function in \eqref{eq:cd-NS}, with $\beta=\beta^{\rm part}+c$ given by
\eqref{eq:cd-beta}. Let $\pi$ be the cover map to $B^\circ$, and define
\[
 \mathcal E_c=\pi\!\left(\bigcup_{E\ne0,\ \operatorname{Pf}(E)=0}Z_E(c)\right),
 \qquad
 \mathcal A_c=\pi\!\left(\bigcup_{\operatorname{Pf}(E)<0}Z_E(c)\right).
\]
These give the exact loci of all elliptic quotients and all abelian
surfaces, respectively. The algebraic dimension of every smooth fibre is
\[
 a(F_t)=
 \begin{cases}
 2,&t\in\mathcal A_c,\\
 1,&t\in\mathcal E_c\setminus\mathcal A_c,\\
 0,&t\notin\mathcal E_c\cup\mathcal A_c.
 \end{cases}
\]
Together with \eqref{eq:cd-NS}, these are an exact classification for every
$t$, by countable unions of explicitly given analytic zero sets; they do
not assert that the loci have a finite algebraic description. They include
all quotient maps, without requiring preservation of the first coordinate.
\end{theorem}
\begin{proof}
We apply the Appell--Humbert correspondence and the Riemann projectivity
criterion in \cite[Theorems 4.1 and 4.4]{Schnell19}; the reduction to the
present explicit periods is as follows. The real vector-space isomorphism
$\Pi:\mathbb R^4\to\mathbb C^2$ transfers $E$ to the actual alternating
form $\widetilde E(v,w)=E(\Pi^{-1}v,\Pi^{-1}w)$ on the tangent space.
Equation \eqref{eq:cd-NS} gives $J$-invariance. The explicit comparison is
\[
 H_E(v,w)=\widetilde E(v,Jw)-i\widetilde E(v,w).
\]
Direct substitution gives $H_E(Jv,w)=iH_E(v,w)$ and
$H_E(w,v)=\overline{H_E(v,w)}$. Its associated real Chern form is
$-\operatorname{Im}H_E=\widetilde E$, since
$(i/2)(H_E(v,w)-H_E(w,v))=\widetilde E(v,w)$.
Thus the Appell--Humbert integrality condition is exactly integrality
of the given $E$, with its sign fully specified. The complex-oriented square computed in D.1.4 shows that
$\operatorname{Pf}(E)<0$ is equivalent to definiteness of this Hermitian
form. Its sign can be chosen positive. The cited projectivity criterion
then gives precisely the asserted abelian locus, in both directions.

For nonzero $E$ with Pfaffian zero, its real rank is two. Its nullspace
$K_{\mathbb R}\subset\mathbb R^4$ is a rational plane because it is the
kernel of an integral matrix. $J$-invariance makes $\Pi K_{\mathbb R}$ a
complex line, and $K=K_{\mathbb R}\cap\mathbb Z^4$ is saturated. The maps
\[
 (\Pi K_{\mathbb R})/\Pi K\ \hookrightarrow F_t
 \ \longrightarrow\
 \mathbb C^2/(\Pi K_{\mathbb R}+\Pi\mathbb Z^4)
\]
are holomorphic group morphisms; the first is a closed embedded elliptic
curve and the second is its quotient onto an elliptic curve, with
connected kernel. A basis of the saturated $K$ extends to an integral basis of $\mathbb Z^4$.
The other two periods project to real-linearly independent vectors in the
one-dimensional complex quotient, because $\Pi$ is a real isomorphism.
They therefore form a discrete lattice there. Compactness follows by
taking the image of its compact real fundamental parallelogram. Conversely any elliptic quotient
pulls an integral positive area form back to a nonzero rank-two integral
$(1,1)$ form. This proves the characterization of $\mathcal E_c$.
Equivalently all its kernels can be found by saturated rank-two lattices
with a basis $v,w$ satisfying
$\det_{\mathbb C}(\Pi v,\Pi w)=0$ and real independence. This is a second
exact integer test, with explicit inclusion and quotient morphisms.

For the algebraic-dimension assertion, an effective divisor on a torus
has a semipositive Hermitian form: averaging its integration current under
translations gives that constant form. A nonzero semipositive form has
rank one or two. Rank two gives the abelian case. Rank one gives the
elliptic quotient just constructed. In the absence of a definite form,
there cannot be two different such kernel subtori. Indeed addition from
their product onto $F_t$ would be a surjective holomorphic homomorphism
with finite kernel, and the product polarization would descend after
multiplication by an integer, giving a definite form. Thus there is a
single kernel $K_t$ in this case. Every effective divisor is a sum of
its translates: its intersection number with $K_t$ is zero, whereas a
curve mapping nonconstantly to the elliptic quotient has strictly positive
intersection with a general quotient fibre. Every meromorphic function
therefore factors through this quotient. One can check the last statement
on a fibre avoiding its polar divisor, where it is a holomorphic function
on a compact elliptic curve and hence constant; analytic continuation
then gives the factorization. The quotient supplies nonconstant
meromorphic functions, so the algebraic dimension is exactly one.
If neither locus occurs, a nonconstant meromorphic function would have
a polar divisor of one of the two forbidden ranks. This proves dimension
zero. Finally monodromy acts by integral determinant-one changes of basis,
preserving \eqref{eq:cd-NS} and its Pfaffian. Thus the stated unions descend
without any marking ambiguity. Each individual matrix and every locus
is specified; no expected identification is left in these formulas.
\end{proof}

\begin{proposition}[D.1.7; Proved]\label{st:D.1.7}
There is an explicit nonempty locus of admissible members with ordinary
fibres of algebraic dimension one. On a sufficiently small punctured
cusp neighbourhood set
\[
 c(z)=-6\tau(z)-12\mu(z)-\beta^{\rm part}(z).
\]
Outside a countable set of $z$, the fibre of $X_{c(z)}$ over $\pi(z)$ has
Neron--Severi lattice $\mathbb Z\eta+\mathbb ZE$, where the nonzero upper
entries of $E$ are
$E_{13}=-6$, $E_{14}=-6$, $E_{23}=1$, $E_{24}=1$.
Its algebraic dimension is one. The section $P(\pi(z))$ can simultaneously
be required to be non-torsion.
\end{proposition}
\begin{proof}
At the specified parameter $\beta=-6\tau-12\mu$. Both $E$ and $\eta$
satisfy \eqref{eq:cd-NS}, and exact multiplication gives
$\operatorname{Pf}(a\eta+bE)=6a^2$.
For any further integral class $F$, substitution in its period equation
gives a linear relation among $1,\tau,\mu,(\mu+\tau)^2$.
These four functions are linearly independent: applying $g_0$ to a
relation $A+B\tau+C\mu+D(\mu+\tau)^2=0$ and subtracting yields
$-B+D(1-2(\mu+\tau))=0$. A further cusp difference forces $D=0$,
and then $B=0$; nonconstancy of $\mu$ forces $A=C=0$.
Comparison of coefficients now gives
$F_{34}=F_{12}=0$, $F_{13}=-6F_{24}$, and
$6F_{23}-F_{14}=12F_{24}$, exactly
$F\in\mathbb Z\eta+\mathbb ZE$. Every other integral class therefore
has a discrete zero set along this period locus. Excluding their
countable union gives the asserted exact Neron--Severi lattice.
It contains the nonzero isotropic $E$ and no matrix of negative
Pfaffian; D.1.6 gives algebraic dimension one.
Finally
$c=-5\tau-12\mu-(\beta^{\rm part}+\tau)$ has imaginary part tending
to minus infinity at the cusp. A sufficiently small cusp neighbourhood
is therefore admissible by C.2.1. The torsion locus of $P$ is countable
by D.1.2, so its exclusion still leaves such points in every open subset
of that neighbourhood. This supplies actual members and fibres, with
all their parameter functions specified.
\end{proof}

\subsection{D.2. Compact curves and surfaces}
\begin{proposition}[D.2.1; Proved]\label{st:D.2.1}
For a standard-twist member \(X_c\) with underlying manifold \(S^6\),
every compact irreducible complex curve lies in one fibre of \(f\), and
every compact irreducible complex surface is one of the reduced fibres.
\end{proposition}
\begin{proof}
The given diffeomorphism with \(S^6\) implies \(H^2(X_c;\mathbb Z)=0\).
For a compact irreducible curve \(C\), let
\(\nu:\widetilde C\to C\hookrightarrow X_c\) be its normalization map.
The holomorphic map \(f\nu\) has
\(\deg (f\nu)^*\mathcal O_{\mathbb P^1}(1)=0\), since its Chern class is
pulled back from the zero class on \(X_c\). A nonconstant holomorphic map
of compact Riemann surfaces has positive degree: a regular value has a
positive number of preimages, all with positive local multiplicity.
Therefore \(f\nu\), and hence \(f|_C\), is constant.

An irreducible surface \(Z\) in the smooth threefold is an effective
Cartier divisor. If its image were all of \(B\), a general smooth fibre
would meet it in a nonempty effective divisor \(D\) whose line bundle
is \(\mathcal O_{X_c}(Z)|_{F_t}\). Its integral Chern class is zero.
Choose a flat K\"ahler form \(\omega\) on the torus \(F_t\). The equality
\(\int_{F_t}c_1(D)\wedge\omega=\int_D\omega\) has left side zero and right
side positive, a contradiction. Thus \(Z\) lies in a single fibre.
Each reduced fibre is irreducible of dimension two, so containment forces
equality. This uses the full map into \(X_c\), not a transfer of a
vanishing claim from its normalization.
\end{proof}

\begin{proposition}[D.2.2; Proved]\label{st:D.2.2}
There is no rational curve in \(X_c\setminus W_c\). There are elliptic
curves outside all three special fibres, in particular in each of the
two smooth fibres \(F_{t_\pm}\). All curves in \(W_c\) are precisely the
images of irreducible algebraic curves in \(\mathrm{dP}_6\) under its
finite normalization map. Distinct curves have the same image only for
the three paired boundary lines. All curves in a bielliptic fibre
\(S_j\) are images of irreducible curves in its finite product cover.
\end{proposition}
\begin{proof}
By D.2.1 every curve is vertical. A holomorphic map from
\(\mathbb P^1\) into a torus lifts to its vector-space universal cover,
since \(\mathbb P^1\) is simply connected, and its coordinate functions
are constant by compactness. Thus a smooth torus fibre contains no
rational curve. A map from \(\mathbb P^1\) into a bielliptic surface
lifts through its finite unramified torus cover by the same covering
criterion; the lifted map is constant. These arguments prove the first
assertion, including possibly singular rational images.
D.1.1 gives an embedded connected elliptic kernel in each \(F_{t_\pm}\),
proving the second.

For the third assertion let \(\eta:\mathrm{dP}_6\to W_c\) be the finite
normalization morphism. Each component of the inverse image of a compact
curve is compact and one-dimensional and maps onto that curve, because
finite maps have zero-dimensional fibres. It is algebraic in the
projective surface \(\mathrm{dP}_6\). Conversely its image is a compact
analytic curve by properness. Away from the conductor, \(\eta\) is a
biholomorphism, so two distinct curves can have the same image only when
both lie in the conductor; its identifications are exactly the three
opposite pairs computed in C.4.1. This is also an explicit algebraic
parametrization: the curves of \(\mathrm{dP}_6\) are the three exceptional
lines and strict transforms of irreducible plane curves under the blow-up
of the three coordinate points. Their homogeneous equations and their
three multiplicities specify the strict transforms. The finite-product
cover statement for \(S_j\) follows from the identical proper-image and
finite-preimage argument applied to its unramified morphism
\(E_j\times F_j\to S_j\); curves in the product are algebraic.
\end{proof}

\begin{theorem}[D.2.3; Known, with derivation]\label{st:D.2.3}
The following descriptions give all compact irreducible curves of $X_c$.
In $W_c$ and $S_j$ they are exactly the images described in D.2.2.
In a smooth fibre of algebraic dimension zero there are none. In a
smooth fibre of algebraic dimension one they are exactly the translates
of its unique elliptic kernel $K_t$. In an abelian fibre they are exactly
the translates of elliptic subtori, together with the irreducible zero
divisors of the positive Appell--Humbert theta sections described below.
These descriptions, D.1.6, and D.2.1 constitute a classification including
every smooth exceptional fibre and every compact irreducible surface.
\end{theorem}
\begin{proof}
All curves are vertical by D.2.1, so the list of fibre types is exhaustive.
The dimension-zero and dimension-one assertions, including uniqueness,
were proved in D.1.6. Let now $F_t=\mathbb C^2/\Pi\Lambda$ be abelian.
For every positive definite Hermitian form $H$ whose imaginary part takes
integer values on $\Pi\Lambda$, and every function
$\alpha:\Pi\Lambda\to U(1)$ satisfying
\[
 \alpha(\lambda+\nu)=\alpha(\lambda)\alpha(\nu)
                 e^{\pi i\operatorname{Im}H(\lambda,\nu)},
\]
take the nonzero entire solutions $\theta:\mathbb C^2\to\mathbb C$ of
\begin{equation}\label{eq:cd-theta-sections}
 \theta(\zeta+\lambda)=\alpha(\lambda)
 e^{\pi H(\zeta,\lambda)+\pi H(\lambda,\lambda)/2}\theta(\zeta).
\end{equation}
The cocycle identity follows by multiplying the displayed factors; its
integral imaginary cross-term has exponential one. It defines a
holomorphic line bundle by quotienting $\mathbb C^2\times\mathbb C$.
The invariant zero set of $\theta$ descends to a compact divisor on
$F_t$. Retain precisely those zero divisors that are reduced and
irreducible. This is an exact parametrization, with all periods and
multiplier equations given; multiplying $\theta$ by a nonzero scalar
does not change its divisor.

The correspondence \cite[Theorem 4.1 and equation (4.3)]{Schnell19} proves
that every holomorphic line bundle and every section on this actual
quotient has this form. For the cohomology matrix the exact comparison is
$H_E=\widetilde E(\,\cdot\,,J\,\cdot\,)-i\widetilde E$ from D.1.6;
we choose the sign of $E$ making $H_E$ positive. Its imaginary part is
$-\widetilde E$, which gives the integral semicharacter equation displayed
above. A divisor class with positive Hermitian rank two gives exactly
these positive theta sections. A divisor of Hermitian rank one is pulled
back from its elliptic quotient and, when irreducible and reduced, is a
single fibre. Thus it is a translate of an elliptic subtorus. The null
lattice is computed by the explicit integer kernel in D.1.6. No other
rank is possible for an effective nonempty divisor, by positivity of its
averaged current. This proves exhaustiveness on the abelian fibre.

The same equations give the curves in the product covers of $S_j$:
compute them on its two-elliptic-curve lattice and take images under the
specified finite morphism. Conversely each curve downstairs is the image
of a component of its finite inverse image, so that this procedure omits
none. In $W_c$ the analogous procedure uses the homogeneous polynomial
parametrization of the normalization in D.2.2, with the three boundary
pairings already identified. The maps used throughout are finite or
closed embeddings, so they preserve compactness and give one-dimensional
images; no assertion is transferred from a normalization in an invalid
direction.
\end{proof}
